\section{Explorando o Sistema}

\begin{frame}{Simplificação: Volumes Constantes}
  \begin{itemize}
    \item Os volumes são variáveis integradoras.
    \item Isto é, qualquer desbalanço entre entrada e saída faz o volume crescer ou diminuir infinitamente.
    \item Para esta análise, assumimos \textbf{volumes constantes} e eliminamos esses integradores.
  \end{itemize}
  Assim, as vazões de saída passam a ser determinadas por relações algébricas:
  \begin{flalign*}
    F_1 &= F_{f1} + F_R \\
    F_2 &= F_{f2} + F_1 \\
    F_3 &= F_2 - F_P - F_R
  \end{flalign*}
  Com isso, fazemos a redução do modelo de \textbf{12} para \textbf{9 estados}.
\end{frame}

\begin{frame}{Ponto Estacionário}
  \begin{columns}
    \column{0.5\textwidth}
    \begin{table}
      \caption{Entradas $u_{ss}$}
      \begin{tabular}{lll}
        \hline
        Variável & Valor & Unidade \\
        \hline
        $F_{f1}$ & 5.04   & m³/h \\
        $F_{f2}$ & 5.04   & m³/h \\
        $F_R$    & 17.00  & m³/h \\
        $Q_1$    & 715.3  & MJ/h \\
        $Q_2$    & 579.8  & MJ/h \\
        $Q_3$    & 568.7  & MJ/h \\
        $T_0$    & 359.1  & K    \\
        \hline
      \end{tabular}
    \end{table}
    \column{0.5\textwidth}
    \begin{table}
      \caption{Estados $y_{ss}$}
      \begin{tabular}{lll}
        \hline
        Variável  & Valor & Unidade \\
        \hline
        $T_1$     & 432.4 & K  \\
        $T_2$     & 427.1 & K  \\
        $T_3$     & 432.1 & K  \\
        $x_{A1}$  & 0.536 & -- \\
        $x_{B1}$  & 0.448 & -- \\
        $x_{A2}$  & 0.545 & -- \\
        $x_{B2}$  & 0.438 & -- \\
        $x_{A3}$  & 0.298 & -- \\
        $x_{B3}$  & 0.670 & -- \\
        \hline
      \end{tabular}
    \end{table}
  \end{columns}
\end{frame}

\begin{frame}{Resposta ao degrau em $F_R$}
  \begin{figure}
    \centering
    \includegraphics[width=0.8\linewidth]{figuras/step-FR-inputs.png}
    \includegraphics[width=0.8\linewidth]{figuras/step-FR-outputs.png}
    % \only<2>{\includegraphics[width=0.8\linewidth]{figuras/step-FR-outputs-xc.png}}
    \caption{Entradas e saídas do sistema para um degrau em $F_R$.}
  \end{figure}
\end{frame}

\begin{frame}
  \newcommand{\ua}{$\uparrow$}
  \newcommand{\da}{$\downarrow$}
  \newcommand{\nz}{$\approx 0$}
  \newcommand{\nm}{$^{\text{nm}}$}
  \begin{columns}
    \column{0.6\textwidth}
    \small
    \begin{table}[h!]
      \centering
      \begin{tabular}{c|ccccccc}
        \hline
        $\mathbf{y} / \mathbf{u}$ & $F_{f1}$ & $F_{f2}$ & \textcolor{red}{$F_R$} & $Q_1$ & $Q_2$ & $Q_3$ & $T_0$ \\
        \hline
        $T_1$    & \da & \da & \textcolor{red}{\da\ua} & \ua    & \ua    & \ua    & \ua    \\
        $x_{A1}$ & \ua & \ua & \textcolor{red}{\ua\da} & \da    & \da    & \da    & \da    \\
        $x_{B1}$ & \da & \da & \textcolor{red}{\da\ua} & \ua    & \ua    & \ua    & \ua\ua \\
        % $x_{C1}$ & \da & \da & \textcolor{red}{\da\da\nz & \ua\nz & \ua\nz & \ua\nz & \ua    \\
        \hline
        $T_2$    & \da & \da & \textcolor{red}{\ua}    & \ua    & \ua    & \ua    & \ua    \\
        $x_{A2}$ & \ua & \ua & \textcolor{red}{\ua\da} & \da    & \da    & \da    & \da    \\
        $x_{B2}$ & \da & \da & \textcolor{red}{\da\ua} & \ua    & \ua    & \ua    & \ua\ua \\
        % $x_{C2}$ & \da & \da & \textcolor{red}{\da\da\nz} & \ua\nz & \ua\nz & \ua\nz & \ua    \\
        \hline
        $T_3$    & \da & \da & \textcolor{red}{\da\ua}    & \ua    & \ua    & \ua    & \ua    \\
        $x_{A3}$ & \ua & \ua & \textcolor{red}{\da}       & \da    & \da    & \da    & \da    \\
        $x_{B3}$ & \da & \da & \textcolor{red}{\ua}       & \ua    & \ua    & \ua    & \ua\da \\
        % $x_{C3}$ & \da & \da & \textcolor{red}{\ua\nm\nz} & \ua\nz & \ua\nz & \ua\nz & \ua    \\
        \hline
      \end{tabular}
      \caption{Interação entre entradas e estados.}
    \end{table}
    \column{0.4\textwidth}
    \begin{itemize}
      \item \ua: ganho positivo;
      \item \da: ganho negativo;
      \item \da\ua: ganho positivo com resposta inversa;
      \item \ua\da: ganho negativo com resposta inversa;
      \item \ua\ua: ganho positivo com sobressinal;
      \item \da\da: ganho negativo com sobressinal;
        % \item \nm: resposta não monotônica;
        % \item \nz: ganho muito pequeno.
    \end{itemize}
  \end{columns}
\end{frame}
